% Standalone source: compile this file directly when a Chapter 5-only PDF is needed.
% Embedded: define \CGAChapterFiveEmbedded before \input{sections/05_theory}.
% In embedded mode the host provides theorem environments and bibliography.
\ifdefined\CGAChapterFiveEmbedded
\else
\documentclass[11pt]{article}
\usepackage[a4paper,margin=25mm]{geometry}
\usepackage{amsmath,amssymb,amsthm,mathtools,booktabs,microtype,enumitem}
\usepackage[hidelinks]{hyperref}
\def\CGAChapterFiveStandalone{1}
\newtheorem{theorem}{Theorem}[section]
\newtheorem{lemma}[theorem]{Lemma}
\newtheorem{proposition}[theorem]{Proposition}
\newtheorem{corollary}[theorem]{Corollary}
\newtheorem{assumption}[theorem]{Assumption}
\theoremstyle{remark}
\newtheorem{remark}[theorem]{Remark}
\setlength{\parskip}{3pt}
\emergencystretch=2em
\begin{document}
\setcounter{section}{4}
\fi
\providecommand{\cE}{\mathcal E}
\providecommand{\cD}{\mathcal D}
\providecommand{\cP}{\mathcal P}
\providecommand{\cK}{\mathcal K}
\providecommand{\cnorm}[1]{\left\lVert#1\right\rVert}
\providecommand{\cabs}[1]{\left|#1\right|}
\providecommand{\cinner}[2]{\left\langle#1,#2\right\rangle}
\providecommand{\cdd}{\,\mathrm d}
\providecommand{\cR}{\mathbb R}

\ifdefined\CGAChapterFiveStandalone
  \newcommand{\CGAHLDreference}{Section~3}
  \newcommand{\CGAEnergyMetricReference}{the energy--distance theorem and metric corollary in Section~4}
  \newcommand{\CGAEnergyTheoremReference}{the energy--distance theorem in Section~4}
\else
  \newcommand{\CGAHLDreference}{Section~\ref{sec:descent}}
  \newcommand{\CGAEnergyMetricReference}{Theorem~\ref{thm:energy-v} and Corollary~\ref{cor:metric-rates}}
  \newcommand{\CGAEnergyTheoremReference}{Theorem~\ref{thm:energy-v}}
\fi

\section{Conditional convergence and finite-trajectory projection bounds}
\label{sec:c5}

We distinguish descent for the defined nonlinear method from approximation
in a frozen Hilbert metric. The first involves continuous-dictionary coverage,
score accuracy, and coefficient correction. The second concerns the nested
spaces generated by the algorithm and does not require the nonlinear method to
maximize a frozen score. The two are connected only conditionally, and for the
quartic energies the connection can also be expressed by an exact identity.

Throughout, $m$ counts accepted independent directions. A rejected attempt or
a direction already in the active span does not increment this index.
The starting index is $m_0$ and $j=m-m_0$. We use $\widehat G_m$ for a computed
state, $G_m^\sharp$ for an exact active-space energy minimizer, and $z_m$
for a frozen orthogonal projection. These states need not coincide.
All derivatives below act on the specified feasible linear space, so that
$\cE'(u^*)=0$ there. Affine constraints are reduced to this setting by a
fixed feasible lift.

\subsection{Finite-pool descent and the residual-comparison condition}
\label{sec:c5-descent}

Let $X$ be a real Banach space, $\cE:X\to\cR$ a differentiable convex
functional on the states and segments used below, and $u^*$ a minimizer.
The symmetric dictionary $\cD_X$ consists of $X$-unit vectors. Set
\[
 \widehat\delta_m=\cE(\widehat G_m)-\cE(u^*),\qquad
 S_m^X=\sup_{g\in\cD_X}\cabs{\cinner{\cE'(\widehat G_m)}g}.
\]
The exact finite-pool score is $s_m(g)=|\langle\cE'(\widehat G_m),g\rangle|$.
The available computed scores satisfy
\begin{equation}
 \cabs{\widetilde s_m(g)-s_m(g)}\leq\xi_m\quad(g\in\cP_m),\qquad
 \widetilde s_m(g_{m+1})\geq t_m^{\rm pool}
                  \max_{g\in\cP_m}\widetilde s_m(g),
 \label{eq:c5-oracle}
\end{equation}
where $\cP_m\subset\cD_X$ is finite and $0<t_m^{\rm pool}\leq1$.
When $S_m^X>0$, define its continuous coverage by
\begin{equation}
 \tau_m=\frac{\max_{g\in\cP_m}s_m(g)}{S_m^X}\in[0,1].
 \label{eq:c5-coverage}
\end{equation}
Set $\tau_m=1$ when $S_m^X=0$. The selection weakness $t_m^{\rm pool}$
and the coverage $\tau_m$ are different quantities.

Assume that a common $L>0$ satisfies
\begin{equation}
 \cE(\widehat G_m+\lambda g)
 \leq\cE(\widehat G_m)+\lambda\langle\cE'(\widehat G_m),g\rangle
            +\tfrac L2\lambda^2,\qquad |\lambda|\leq1,\quad g\in\cD_X.
 \label{eq:c5-directional}
\end{equation}
The expanded space contains $\widehat G_m+\lambda g_{m+1}$, and the
correction satisfies
\begin{equation}
 \cE(\widehat G_{m+1})\leq
 \inf_{v\in\widehat\Phi_{m+1}}\cE(v)+\varepsilon_m^{\rm opt},
 \qquad \varepsilon_m^{\rm opt}\geq0 .
 \label{eq:c5-correction}
\end{equation}
This assumption concerns the true energy. A numerical objective needs its
own comparison with that energy.

\begin{proposition}[Descent with a resolved score]
\label{prop:c5-descent}
Let $a_m\geq0$ be any valid lower bound for $s_m(g_{m+1})$.
Under \eqref{eq:c5-directional}--\eqref{eq:c5-correction},
\begin{equation}
 \widehat\delta_{m+1}\leq\widehat\delta_m-H_L(a_m)
                       +\varepsilon_m^{\rm opt}.
 \label{eq:c5-descent}
\end{equation}
Here $H_L$ is the directional-decrease function defined in \CGAHLDreference.
In particular, \eqref{eq:c5-oracle} gives the valid lower bound
\begin{equation}
 a_m=\left[t_m^{\rm pool}\tau_m S_m^X
                 -(1+t_m^{\rm pool})\xi_m\right]_+ .
 \label{eq:c5-score-lower}
\end{equation}
\end{proposition}
\begin{proof}
The triangle inequality gives
$s_m(g_{m+1})\geq
t_m^{\rm pool}\max_{\cP_m}s_m-(1+t_m^{\rm pool})\xi_m$.
Test the correction with a signed dictionary step in the true descent
direction of length $\min\{1,a_m/L\}$. Both signs lie in the expanded
linear space. Minimizing the quadratic upper bound on that segment
gives $H_L(a_m)$ and proves \eqref{eq:c5-descent}.
\end{proof}

The radial stationarity error enters when an atomic source is used to
bound $S_m^X$. Write $\cK_1(\cD)$ for the closed absolutely convex hull
of a bounded dictionary. If
$u^*\in B_{\rm src}\cK_1(\cD_X)$ with $B_{\rm src}>0$ and
$|\langle\cE'(\widehat G_m),\widehat G_m\rangle|
\leq\eta_m^{\rm rad}$, convexity implies
\begin{equation}
 \widehat\delta_m\leq\eta_m^{\rm rad}+B_{\rm src}S_m^X .
 \label{eq:c5-atomic}
\end{equation}
Consequently one may instead use
\[
 a_m=\left[
 t_m^{\rm pool}\tau_m
       \frac{\widehat\delta_m-\eta_m^{\rm rad}}{B_{\rm src}}
 -(1+t_m^{\rm pool})\xi_m\right]_+ .
\]
This is a source-dependent estimate, and it is not an additional error
term to be added after the score estimate \eqref{eq:c5-score-lower}.
In the error framework of \cite[Section~3, Theorem~3.2 and
Lemma~6.3]{DereventsovTemlyakov2022}, energy correction and radial
biorthogonality are separate errors. Their power-smoothness rate assumes
uniform positive continuous weakness and prescribed decay of those errors.
Here they correspond to $\varepsilon_m^{\rm opt}$ and
$\eta_m^{\rm rad}$, with the latter at the preceding state in
\eqref{eq:c5-atomic}. Pool coverage and score quadrature determine a further
effective weakness through \eqref{eq:c5-score-lower}. Accuracy of the
computed correction alone does not ensure continuous weakness.

\begin{lemma}[A nonautonomous scalar bound]
\label{lem:c5-recurrence}
Suppose $x_m\geq0$ and
$x_{m+1}\leq x_m-b_mx_m^q$ for $q>1$ and $b_m\geq0$.
For $x_0>0$,
\begin{equation}
 x_n\leq\left[x_0^{1-q}+(q-1)\sum_{m=0}^{n-1}b_m\right]^{-1/(q-1)} .
 \label{eq:c5-recurrence}
\end{equation}
If $x_0=0$, the sequence vanishes. In particular,
$b_m\geq b(m+1)^{-s}$, $b>0$, $0\leq s<1$, implies
$x_n=O((n+1)^{-(1-s)/(q-1)})$.
\end{lemma}
\begin{proof}
As long as $x_{m+1}>0$, convexity gives
$x_{m+1}^{1-q}-x_m^{1-q}\geq
(q-1)x_m^{-q}(x_m-x_{m+1})\geq(q-1)b_m$.
Sum. Once a term vanishes, nonnegativity and the recurrence force all
later terms to vanish. The power estimate follows by comparison of the
sum with an integral.
\end{proof}

\begin{theorem}[Conditional residual-comparison rate]
\label{thm:c5-visible}
Let $p>2$. Suppose, at the computed states,
\begin{align}
 \widehat\delta_m&\geq c_U\cnorm{\widehat G_m-u^*}_X^p,\quad c_U>0,
 \label{eq:c5-uniform-convex}\\
 S_m^X&\geq\kappa_0(m+1)^{-\chi}
                   \cnorm{\cE'(\widehat G_m)}_{X^*},
 \quad\kappa_0>0,\quad0\leq\chi<\tfrac12 .
 \label{eq:c5-visible}
\end{align}
Assume \eqref{eq:c5-oracle}--\eqref{eq:c5-correction}, and, uniformly in $m$,
\[
 t_m^{\rm pool}\tau_m\geq\underline t>0,\qquad
 (1+t_m^{\rm pool})\xi_m
 \leq\vartheta t_m^{\rm pool}\tau_m S_m^X,\quad0\leq\vartheta<1 .
\]
With $a_m$ as in \eqref{eq:c5-score-lower}, assume also
$\varepsilon_m^{\rm opt}\leq\zeta H_L(a_m)$, $0\leq\zeta<1$.
Then
\begin{equation}
 \widehat\delta_m=O((m+1)^{-p(1-2\chi)/(p-2)}) .
 \label{eq:c5-visible-rate}
\end{equation}
The constant depends on $p,c_U,\kappa_0,\underline t,L,\vartheta,\zeta$ and
$\widehat\delta_0$. For the exact continuous method take
$\tau_m=1$, $\xi_m=\varepsilon_m^{\rm opt}=0$.
\end{theorem}
\begin{proof}
Convexity and \eqref{eq:c5-uniform-convex} give
\[
 \widehat\delta_m\leq
 \cnorm{\cE'(\widehat G_m)}_{X^*}
       (\widehat\delta_m/c_U)^{1/p}.
\]
For positive error this yields gradient domination.
Thus $a_m\geq b_m^*:=K(m+1)^{-\chi}
\widehat\delta_m^{(p-1)/p}$, where
$K=(1-\vartheta)\underline t\kappa_0c_U^{1/p}$.
Descent and absorption imply monotonicity, so
$b_m^*\leq A:=K\widehat\delta_0^{(p-1)/p}$.
For $0\leq b\leq A$,
$H_L(b)\geq b^2/(2\max\{L,A\})$.
Since $H_L$ is increasing, \eqref{eq:c5-descent} implies
\[
 \widehat\delta_{m+1}\leq\widehat\delta_m-
 \frac{(1-\zeta)K^2}{2\max\{L,A\}}
 (m+1)^{-2\chi}\widehat\delta_m^{2(p-1)/p}.
\]
Apply Lemma~\ref{lem:c5-recurrence}. The case $\widehat\delta_0=0$
is immediate from descent and absorption.
\end{proof}

This theorem does not assert that density of $\operatorname{span}\cD_X$
implies \eqref{eq:c5-visible}. For relatively compact dictionaries in an
infinite-dimensional space, a uniform norm comparison on the entire dual
space is generally impossible. The comparison is therefore a hypothesis on
the specified residual trajectory. The resulting estimate improves on the
atomic energy baseline $m^{-1}$ exactly when $\chi<1/p$.

\subsection{Energy and the natural distance}
\label{sec:c5-natural}

For the standard pure, regularized, and reaction $p$-energies, the energy--distance
equivalence and the $W^{1,p}$ conversion are established in
\CGAEnergyMetricReference. Therefore, if the
conditional energy estimate in Theorem~\ref{thm:c5-visible} is
$\widehat\delta_m=O(m^{-\alpha})$, then
\[
 d_{V,\bullet}(\widehat G_m,u^*)=O(m^{-\alpha/2}),\qquad
 \|\widehat G_m-u^*\|_{W^{1,p}}=O(m^{-\alpha/p}),
\]
under the corresponding zero-mean or reaction assumptions. An upper energy bound
gives an upper metric bound; it is not an equality with a fitted slope.

Thus \eqref{eq:c5-visible-rate} has natural-distance exponent
$p(1-2\chi)/(2(p-2))$ and Sobolev exponent $(1-2\chi)/(p-2)$ for the
pure, regularized, and reaction energies. The scalar comparison in Theorem~\ref{thm:c5-visible}
does not automatically extend to other energy densities.

\subsection{Local Hessian transfer and the same-dictionary requirement}
\label{sec:c5-hessian}

Previous theoretical work provides substantial support for the transfer argument.
Xu--Xu \cite[Theorem~3.1, equations~(17)--(19)]{XuXu2025}
analyze exact active-space minimization of a convex, $L$-smooth
functional in a Sobolev Hilbert space. In their notation, the continuous
selection strengths $\gamma_i$, source radius $B$, and innovations
$\bar g_i=(I-P_{i-1})g_i$ yield
\[
 \delta_i\leq\delta_{i-1}
 -\frac{\gamma_i^2\delta_{i-1}^2}{2LB^2\cnorm{\bar g_i}^2}.
\]
Their ReLU dictionary estimate then gives an energy rate
$n^{-1-(2(k-s)+1)/d}$ for constant positive weakness in $H^s$,
with $k\geq s$; strong convexity gives the corresponding norm bound.
Here $s$ is the Sobolev order, not the iteration index.
That exponent belongs to their normalized dictionary and hypotheses; it is
not imported into the present finite, discretely normalized pool without a
separate compatible entropy proof.
The argument below states explicit sufficient conditions for
transferring nonlinear selection to a frozen metric. It does not establish
a new general nonlinear greedy rate. Global smoothness, exact correction,
and the dictionary assumptions in the direct theorem cannot be inferred
from a finite-pool computation or from local Hessian positivity.

Relative to Dereventsov--Temlyakov, the finite-pool proposition below
separates coverage, score-quadrature, radial-stationarity, and
coefficient-correction errors. Relative to Xu--Xu, the transfer theorem
retains their exact active-space and entropy mechanism and adds only a
conditional local Hessian transfer. Relative to Li--Siegel, the entropy
lemma specializes the volume argument to weak selection and unrenormalized
projected atoms. These comparisons identify the decomposition of the selection,
quadrature, stationarity, and correction terms used in the present transfer argument;
they do not by themselves establish a new nonlinear rate.

Let $H$ be a real Hilbert space on which $\cE$ is twice Fr\'echet
differentiable in an open convex neighborhood of $u^*$.
Assume $A_*=\cE''(u^*)$ is bounded, symmetric, and coercive on $H$.
Write $a_*(v,w)=\langle A_*v,w\rangle$ and $\cnorm v_*^2=a_*(v,v)$.
Suppose the closed ball of radius $R_*$ around $u^*$ lies in that
neighborhood and
\begin{equation}
 |\langle[\cE''(u)-A_*]v,w\rangle|
 \leq L_*\cnorm{u-u^*}_*\cnorm v_*\cnorm w_*,
 \qquad \cnorm{u-u^*}_*\leq R_*,\quad L_*R_*\leq\tfrac12 .
 \label{eq:c5-hessian}
\end{equation}
For finite-dimensional $\Phi\subset H$ let $z_\Phi$ be the frozen
orthogonal projection of $u^*$ and $G_\Phi^\sharp$ an exact minimizer
of $\cE$ over $\Phi$. Assume both points lie in this ball.

\begin{lemma}[Taylor and projection estimates]
\label{lem:c5-projection}
Under \eqref{eq:c5-hessian}, for $u$ in the ball,
\begin{align}
 \cabs{\cE(u)-\cE(u^*)-\tfrac12\cnorm{u-u^*}_*^2}
 &\leq\tfrac{L_*}{6}\cnorm{u-u^*}_*^3,\label{eq:c5-taylor}\\
 \cnorm{\cE'(u)-A_*(u-u^*)}_{*'}
 &\leq\tfrac{L_*}{2}\cnorm{u-u^*}_*^2,\label{eq:c5-grad-taylor}\\
 \cnorm{G_\Phi^\sharp-z_\Phi}_*
 &\leq L_*\cnorm{z_\Phi-u^*}_*^2.\label{eq:c5-proj}
\end{align}
Here the dual norm in \eqref{eq:c5-grad-taylor} is taken over $H$.
\end{lemma}
\begin{proof}
One and two integrations of \eqref{eq:c5-hessian} give the first
two estimates. Put $d=G_\Phi^\sharp-z_\Phi\in\Phi$.
Use $\langle\cE'(G_\Phi^\sharp),d\rangle=0$ and
$a_*(z_\Phi-u^*,d)=0$. Subtracting $A_*d$ along the segment from
$z_\Phi$ to $G_\Phi^\sharp$ bounds
$\cnorm d_*^2$ by
$L_*R_*\cnorm d_*^2+(L_*/2)\cnorm{z_\Phi-u^*}_*^2\cnorm d_*$.
Absorb the first term. All segments remain in the ball by convexity.
\end{proof}

Assume $\cD_X\subset H$ and the algorithmic dictionary satisfies the uniform norm comparison
\begin{equation}
 0<c_*\leq\cnorm g_*\leq C_*<\infty\quad(g\in\cD_X),\qquad
 \cD_*=\{g/\cnorm g_*:g\in\cD_X\}.
 \label{eq:c5-bridge}
\end{equation}
Let $\Phi_m$ be nested finite-dimensional spaces with
$\Phi_{m+1}=\Phi_m+\operatorname{span}\{g_{m+1}\}$, where each selected
direction is independent of $\Phi_m$. Denote their frozen orthogonal
projections by $P_m$, and set $z_m=P_mu^*$,
$r_m=\cnorm{u^*-z_m}_*$, and
$G_m^\sharp=G_{\Phi_m}^\sharp$.
At an exact active-space state set
\[
 S_m^*=\sup_{d\in\cD_*}|a_*(z_m-u^*,d)|,\qquad
 \eta_m^H=\sup_{g\in\cD_X}
 |\langle\cE'(G_m^\sharp),g\rangle-a_*(z_m-u^*,g)|.
\]
The preceding lemma yields
$\eta_m^H\leq C_T\cnorm{z_m-u^*}_*^2$, with a constant uniform on the
specified ball and dictionary. This absolute bound alone need not be
small compared with $S_m^*$.

\begin{proposition}[Selection under a relative score condition]
\label{prop:c5-transfer}
Suppose nonlinear selection over $\cD_X$ has strength $t\in(0,1]$.
For $0<t_0<tc_*/C_*$, the condition
\begin{equation}
 (1+t)\eta_m^H\leq(tc_*-t_0C_*)S_m^*
 \label{eq:c5-relative}
\end{equation}
implies
$|a_*(z_m-u^*,g_{m+1}/\cnorm{g_{m+1}}_*)|\geq t_0S_m^*$.
\end{proposition}
\begin{proof}
The selected unnormalized frozen score is at least
$t\sup_{g\in\cD_X}|a_*(z_m-u^*,g)|-(1+t)\eta_m^H$.
The supremum is at least $c_*S_m^*$, and division by the selected
norm costs at most $C_*$. This proves the assertion.
\end{proof}

Fix $m_0$, $Q_0=I-P_{m_0}$ and
\[
 H_0=Q_0H,\qquad f_0=Q_0u^*,\qquad
 \cD_{\rm rst}=Q_0\cD_*\setminus\{0\}.
\]
The projected atoms have norm at most one and are not renormalized.
The entropy numbers use the \(2^{n-1}\)-ball convention fixed in
the variational setting; consequently a cover by \(2^j\) balls in the
volume argument below is represented by \(\varepsilon_{j+1}\), not by a
second definition of the same symbol.

\begin{lemma}[Entropy bound for bounded-atom weak OGA]
\label{lem:c5-entropy}
Let $\cD_0$ be symmetric with atom norms at most one in a Hilbert space,
$K=\cK_1(\cD_0)$ compact, and $f_0\in B_*K$ with $B_*\geq0$.
Let weak OGA of strength $t_0>0$ start from the zero space.
After $j\geq1$ independent extensions its error satisfies
\begin{equation}
 \cnorm{f_0-P_jf_0}\leq
 \frac{B_*}{t_0}\frac{(j!V_j)^{1/j}}{\sqrt j}
                    \varepsilon_{j+1}(K),\qquad
 V_j=\frac{\pi^{j/2}}{\Gamma(j/2+1)} .
 \label{eq:c5-entropy}
\end{equation}
\end{lemma}
\begin{proof}
If $B_*=0$, then $f_0=0$ and the assertion is immediate. Suppose $B_*>0$.
Write $y_i=\cnorm{f_0-P_if_0}^2$ and
$l_i=\cnorm{(I-P_{i-1})d_i}>0$.
Orthogonality and the source condition give the best score at least
$y_{i-1}/B_*$. Hence
$y_i\leq y_{i-1}-t_0^2y_{i-1}^2/(B_*^2l_i^2)$.
If the error reaches zero the result is immediate. Otherwise taking
reciprocals and using the arithmetic--geometric mean inequality yields
\[
 y_j\leq \frac{B_*^2}{t_0^2\sum_{i=1}^jl_i^{-2}}
 \leq \frac{B_*^2}{t_0^2j}\Big(\prod_{i=1}^jl_i\Big)^{2/j}.
\]
The cross-polytope generated by $\{\pm d_1,\ldots,\pm d_j\}$ has
$j$-dimensional volume $(2^j/j!)\prod_i l_i$ and lies in the orthogonal
projection of $K$ onto their span. A cover of $K$ by $2^j$ balls of radius
\(\varepsilon_{j+1}(K)+h\), \(h>0\), projects to a cover of that
cross-polytope. Volume comparison gives
$\prod_i l_i\leq j!V_j(\varepsilon_{j+1}(K)+h)^j$.
Let $h\downarrow0$ and substitute. No projected atom was rescaled.
\end{proof}

The argument is the volume/entropy argument of Li--Siegel, specialized here
to weak selection and to atoms bounded by one. It corresponds
to Theorem~4.2 and equation~(4.4) of \cite{LiSiegelEntropy2024}, with the
projection recurrence in equations~(4.5)--(4.9); the numbering is that of the
cited reference. The specialization records the innovation
weights of the restarted spaces; it is not an independent entropy theorem.
The factorial--volume factor divided by $\sqrt j$ is bounded uniformly
in $j$ by Stirling's formula.

\begin{theorem}[Conditional Hessian transfer after restart]
\label{thm:c5-hessian}
Assume \eqref{eq:c5-hessian} and \eqref{eq:c5-bridge}.
For every state in $m_0\leq m\leq m_0+J$, suppose that the exact nonlinear
state and frozen projection lie in the specified ball.
For each selection step $m_0\leq m<m_0+J$, assume nonlinear selection
has strength at least $t$, and \eqref{eq:c5-relative} holds with
the same $t_0>0$. On the unrenormalized restarted dictionary assume
\begin{equation}
 f_0\in B_*\cK_1(\cD_{\rm rst}),\qquad
 \varepsilon_j(\cK_1(\cD_{\rm rst});H_0)\leq C_{\rm ent}j^{-\beta},
 \label{eq:c5-source-entropy}
\end{equation}
for $1\leq j\leq J+1$, with compact closed absolutely convex hull,
$B_*,C_{\rm ent}>0$, and $\beta>0$.
Then
\[
 \cnorm{z_{m_0+j}-u^*}_*\leq
 C B_*C_{\rm ent}t_0^{-1}j^{-\beta},\qquad
 \cE(G_{m_0+j}^\sharp)-\cE(u^*)\leq C_{\rm nl}j^{-2\beta}.
\]
Here $C$ is an absolute constant and $C_{\rm nl}$ additionally depends
on $L_*,R_*,B_*,C_{\rm ent},t_0$. The \(j+1\) index in
Lemma~\ref{lem:c5-entropy} is the harmless consequence of the global
\(2^{n-1}\)-ball convention; the entropy assumption is stated through
\(J+1\) for this reason. For a finite $J$ the conclusion is restricted to
that finite range.
\end{theorem}
\begin{proof}
Proposition~\ref{prop:c5-transfer} gives weak frozen selection.
The residual is orthogonal to the starting space, so its scores against
$d$ and $Q_0d$ agree. The spaces after the starting index therefore form weak OGA
on $H_0$ with dictionary $\cD_{\rm rst}$. Apply
Lemma~\ref{lem:c5-entropy}. Lemma~\ref{lem:c5-projection} gives
$\cnorm{G_{m_0+j}^\sharp-u^*}_*
\leq r_{m_0+j}+L_*r_{m_0+j}^2
\leq(1+L_*R_*)r_{m_0+j}$, and \eqref{eq:c5-taylor}
gives the energy bound.
\end{proof}
Constants uniform in a terminal width are required for an asymptotic
conclusion. No exponent depending on a ReLU degree is substituted
without a proof for this dictionary.

This is a conditional transfer statement. The family below verifies
coercivity, source membership, and compactness, but it does not verify
\eqref{eq:c5-relative} or a quantitative entropy exponent for
\(\lambda>0\). Hence the theorem does not yield an unconditional nonlinear
asymptotic rate for that family.

\paragraph{A family satisfying the structural assumptions.}
Take $H=H^1(0,1)$ and
\[
 \cE_\lambda(u)=\int_0^1
 \left(\tfrac12|u'|^2+\tfrac12u^2+\tfrac\lambda4u^4\right)\cdd x-\ell(u),
 \qquad\lambda\geq0 .
\]
For a fixed integer $k\geq2$ and $0<a<b<1$, put
$g_t=(x-t)_+^k/\cnorm{(x-t)_+^k}_{H^1}$, $t\in[a,b]$,
and $\cD_X=\{\pm g_t\}$.
Choose any finite signed measure $\mu$ with total variation at most
$B_{\rm src}$, set $u^*=\int_a^b g_t\cdd\mu(t)$, and manufacture
$\ell(v)=\int(u^{*\prime}v'+u^*v+\lambda(u^*)^3v)$.
The embedding $H^1(0,1)\hookrightarrow L^\infty$ makes this load
continuous and the Hessian locally Lipschitz on $H^1$ balls.  The load is
defined in variational form; it should not be relabelled as a homogeneous
Neumann load unless the corresponding boundary-flux compatibility has also
been checked.
Its bilinear form is
$\int(v'w'+(1+3\lambda(u^*)^2)vw)$, which is coercive in $H^1$.
The norm comparison \eqref{eq:c5-bridge} follows, the dictionary is compact,
and $u^*\in C_*B_{\rm src}\cK_1(\cD_*)$.
Its closed absolutely convex hull is compact: finite nets for the
dictionary approximate that hull by bounded finite-dimensional convex hulls.
Projection by $Q_0$ preserves the source and cannot increase entropy.
Thus these structural assumptions are simultaneously meaningful for an
infinite family of targets, not only for one atom or one fixed finite space.
For $\lambda>0$, the relative-score condition
\eqref{eq:c5-relative} and a quantitative entropy exponent remain additional
hypotheses. For $\lambda=0$, $\eta_m^H=0$ and the setting reduces to
ordinary Hilbert OGA. This example does not, by itself, establish a new
nonlinear fast rate.

The missing relative-score condition can nevertheless be reduced to a
quantitative small-nonlinearity check.  Let
$C_{\rm emb}=\|H^1(0,1)\hookrightarrow L^\infty(0,1)\|$ (for the
standard squared $H^1$ norm one may use the explicit bound
$C_{\rm emb}\leq\sqrt2$) and
$M_* = \|u^*\|_{L^\infty}$, and use the $H^1$-normalized atoms above.
Since $a_*$ dominates the standard $H^1$ inner product, the quartic
Hessian satisfies, on the specified $R_*$-ball,
\begin{equation}
 L_\lambda:=3\lambda C_{\rm emb}
       (C_{\rm emb}R_*+2M_*),\qquad L_\lambda R_*\leq\tfrac12,
 \label{eq:c5-family-hessian-constant}
\end{equation}
and the norm-comparison and source constants can be chosen as
\begin{equation}
 c_*=1,\qquad C_\lambda=\sqrt{1+3\lambda M_*^2},\qquad
 B_\lambda=C_\lambda B_{\rm src}.
 \label{eq:c5-family-bridge-constants}
\end{equation}
Indeed, $\|g_t\|_*=\sqrt{1+3\lambda\|u^*g_t\|_{L^2}^2}$ lies
between $1$ and $C_\lambda$, and the change of variables from $g_t$ to
$g_t/\|g_t\|_*$ multiplies the atomic budget by at most $C_\lambda$.
Put
\begin{equation}
 C_{T,\lambda}=L_\lambda\left[1+\tfrac12(1+L_\lambda R_*)^2\right].
 \label{eq:c5-family-transfer-constant}
\end{equation}
The projection estimates then give
$\eta_m^H\leq C_\lambda C_{T,\lambda}r_m^2$.  Consequently the
following explicit bound is sufficient for the relative-score condition:
\begin{equation}
 \Delta_\lambda:=(tc_*-t_0C_\lambda)
 -(1+t)C_\lambda^2C_{T,\lambda}B_{\rm src}>0.
 \label{eq:c5-small-lambda-score}
\end{equation}
This is a check on specified constants, not a replacement for the invariant
ball or entropy arguments.  For fixed $R_*$, $M_*$, $B_{\rm src}$, $t$, and
$t_0$, the left-hand nonlinear term tends to zero as $\lambda\downarrow0$;
continuity therefore gives a positive, specified $\lambda_0$ whenever
$t_0<tc_*$, provided the constants are held fixed. The admissible interval
$0<\lambda\leq\lambda_0$ must satisfy both
\eqref{eq:c5-family-hessian-constant} and \eqref{eq:c5-small-lambda-score}
throughout the interval, with the values of $R_*$, $M_*$, $B_{\rm src}$,
$t$, and $t_0$ specified.

\paragraph{Sufficient conditions for the concrete semilinear family.}
The following statement isolates the sufficient conditions needed for a
nonlinear claim. Compactness of the parameter interval alone does not supply any of
their quantitative parts. Fix $m_0$, $Q_0=I-P_{m_0}$, and the unrenormalized restarted
dictionary $\cD_{\rm rst}=Q_0\cD_*\setminus\{0\}$, where
$\cD_*=\{g_t/\|g_t\|_*:t\in[a,b]\}$ with both signs included.  The sufficient conditions on a finite range $0\leq j\leq J$ are:
\begin{enumerate}[label=(A\arabic*)]
  \item an initial radius and a specified $R_*$ for which every exact active
  minimizer, frozen projection, and corrected state in the stated finite range stays in
  $B_*(u^*,R_*)$; for a specified $r_{\rm ent}>0$ this includes
  $r_{m_0}\leq r_{\rm ent}$ and
  $L_\lambda R_*\leq\tfrac12$;
  \item a uniform selection strength $t$ and $t_0<tc_*/C_\lambda$ for
  which the displayed bound $\Delta_\lambda$ is positive at every
  selection step $m_0\leq m<m_0+J$ in the stated range.  This range
  contains no exceptional prefix;
  \item a quantitative entropy estimate for this same projected,
  unrenormalized object, with the ball-counting convention fixed above,
  \begin{equation}
    \varepsilon_n\!\left(\cK_1(\cD_{\rm rst});H_0\right)
       \leq C_{\rm ent}n^{-\beta_A},
       \qquad 1\leq n\leq J+1,
    \label{eq:c5-family-entropy-certificate}
  \end{equation}
  and $Q_0u^*\in B_\lambda\cK_1(\cD_{\rm rst})$; and
  \item correction, quadrature, and radial errors absorbed in the
  finite-pool score bound, with tolerances prescribed from the available finite-pool quantities.
\end{enumerate}
For $k\geq2$, one possible target is $\beta_A=1$, which would follow from a direct
parameter-metric estimate for $t\mapsto g_t$ together with a covering argument for the
absolutely convex hull. Until such an estimate is proved for the same projected dictionary,
only an exponent established by that argument may be used. If the estimate is available only
for a fixed $J$, the resulting conclusion is restricted to that finite range.

\begin{corollary}[Concrete-family conditional rate]
\label{cor:c5-concrete-family}
Assume the sufficient conditions above hold for some $\lambda\in(0,\lambda_0]$
and $0<\beta_A$, and use exact active-space minimization in the statement
of the rate.  Then the restarted CGA trajectory satisfies, for
$1\leq j\leq J$,
\begin{equation}
 \cnorm{z_{m_0+j}-u^*}_*
   \leq C B_\lambda C_{\rm ent}t_0^{-1}j^{-\beta_A},
 \qquad
 \cE_\lambda(G_{m_0+j}^\sharp)-\cE_\lambda(u^*)
   \leq C_{\rm nl}j^{-2\beta_A}.
 \label{eq:c5-concrete-family-rate}
\end{equation}
The constants are those of Theorem~\ref{thm:c5-hessian}; if the conditions
are established only through $J$, \eqref{eq:c5-concrete-family-rate} is not an
asymptotic assertion. An unconditional rate requires the relative-score,
invariant-neighbourhood, and entropy conditions to hold uniformly along the trajectory,
with constants independent of $J$.
\end{corollary}

\begin{remark}[Scope of the conditional family result]
The family construction and the explicit Hessian and norm-comparison formulas reduce
the abstract hypotheses to quantities that can in principle be estimated. They do not
by themselves establish a positive relative-score bound, the entropy estimate in
\eqref{eq:c5-family-entropy-certificate}, or invariance of a computed trajectory.
These remain assumptions to be verified before the conditional theorem can be stated
as an unconditional rate result.
\end{remark}

\paragraph{Scope of the function-space assumptions.}
Full normalization of boundary-concentrating truncated powers need not
give a compact dictionary. For example, as $t\uparrow1$, the function
\[
 b_t(x)=\frac{(x-t)_+^k}{\cnorm{(x-t)_+^k}_{H^1(0,1)}}
\]
has $H^1$ norm one
but $L^2$ norm tending to zero. A strongly convergent $H^1$ subsequence
would therefore have limit zero, a contradiction.
For the two-dimensional sinh model, an $H^1$ ball does not give
uniform $L^\infty$ bounds; a Hessian estimate restricted to bounded
states is not a proof of \eqref{eq:c5-hessian} on an $H^1$ open ball.
Without a separate functional-analytic argument we use that model only
in the finite-dimensional trajectory setting. Degenerate pure $p$-energies
likewise need not satisfy full-space Hessian coercivity.

\subsection{Finite-reference scores and restricted stationarity}
\label{sec:c5-validation}

On a specified finite-dimensional space $Y_M$ with a positive frozen inner
product, assume $u^*\in Y_M$ and $\widehat\Phi_m\subset Y_M$.
In this setting let $z_m$ be the frozen projection of $u^*$ onto
$\widehat\Phi_m$. Let the fixed reference set $\mathcal A_M\subset Y_M$
contain all pools on the stated finite range.
Assume $0<c_M\leq\kappa_g:=\cnorm g_*\leq C_M<\infty$.
Set
\[
 h_m(g)=|a_*(u^*-z_m,g)|,\quad
 S_m^{\rm val}=\max_{g\in\mathcal A_M}h_m(g)/\kappa_g,\quad
 \tau_m^{\rm val}
   =\frac{\max_{g\in\cP_m}h_m(g)/\kappa_g}{S_m^{\rm val}}
\]
when $S_m^{\rm val}>0$.
If the computed scores obey
$|\widetilde s_m(g)-h_m(g)|\leq\Delta_m$ on $\cP_m$, the proof of
Proposition~\ref{prop:c5-transfer} gives
\begin{equation}
 \frac{h_m(g_{m+1})}{\kappa_{g_{m+1}}}
 \geq
 \frac{t_m^{\rm pool}c_M\tau_m^{\rm val}S_m^{\rm val}
             -(1+t_m^{\rm pool})\Delta_m}{C_M}.
 \label{eq:c5-finite-margin}
\end{equation}
Thus the following nonnegative lower bound
$(t_m^{\rm pool}c_M\tau_m^{\rm val}-t_0C_M)S_m^{\rm val}
 -(1+t_m^{\rm pool})\Delta_m$ suffices for weakness $t_0$
relative to this finite reference set.
A negative value of this lower bound yields no conclusion about the actual selected score.
A continuous conclusion additionally requires
$S_m^{\rm val}\geq(1-\zeta_M)S_m^*$ with $0\leq\zeta_M<1$.
Comparing two quadrature rules supplies a comparison of quadrature sensitivity, not an established bound for $\Delta_m$.

The active-space KKT residual in a reference norm $\cnorm\cdot_{0,M}$ is
\begin{equation}
 k_m=\sup_{\substack{v\in\widehat\Phi_m\\\cnorm v_{0,M}=1}}
              |\langle\cE'(\widehat G_m),v\rangle|.
 \label{eq:c5-kkt}
\end{equation}
The supremum is over the active space, even when its norm is inherited
from $Y_M$; set $k_m=0$ if $\widehat\Phi_m=\{0\}$.
If $\cE$ is $\mu_{\rm act}>0$ strongly convex on the active
space and its exact minimizer exists, then
\[
 \cnorm{\widehat G_m-G_m^\sharp}_{0,M}\leq k_m/\mu_{\rm act},\quad
 \cE(\widehat G_m)-\cE(G_m^\sharp)\leq k_m^2/(2\mu_{\rm act}).
\]
The first follows by strong monotonicity paired with the state difference;
the second follows by minimizing the strong-convexity quadratic.
A bound $C_{\rm score}$ on
$\cnorm{\cE''(u)g|_{\widehat\Phi_m}}_{\widehat\Phi_m^*}$ over the
relevant segments and pool gives an additional score error at most
$C_{\rm score}k_m/\mu_{\rm act}$.
All norms and spaces must be specified before using such estimates.
In particular,
\[
 \inf_{0\ne v\in\widehat\Phi_m}\frac{a_*(v,v)}{\cnorm v_{0,M}^2}
 \ \geq\
 \inf_{0\ne v\in Y_M}\frac{a_*(v,v)}{\cnorm v_{0,M}^2}.
\]
A positive active-space minimum does not establish a lower bound on $Y_M$
or on an infinite-dimensional space.

\subsection{An approximate-source estimate for the computed trajectory}
\label{sec:c5-source}

Let $Y$ be a finite-dimensional real Hilbert space containing $u^*$,
all active spaces, and a fixed finite symmetric dictionary
$\cD_{\rm fin}$ of frozen-unit vectors. For $m_0\leq m<m_0+J$ assume
\begin{equation}
 \widehat\Phi_{m+1}=\widehat\Phi_m+\operatorname{span}\{d_{m+1}\},
 \quad d_{m+1}\in\cD_{\rm fin},\quad d_{m+1}\notin\widehat\Phi_m .
 \label{eq:c5-spaces}
\end{equation}
The dictionary must contain every accepted direction in this finite range.
Here $P_m$ denotes the frozen orthogonal projection onto $\widehat\Phi_m$.
Set
\[
 \begin{gathered}
 z_m=P_mu^*,\qquad e_m=u^*-z_m,\qquad r_m=\cnorm{e_m}_*,\\
 S_m^{\rm fin}=\max_{d\in\cD_{\rm fin}}|a_*(e_m,d)|,\qquad
 \ell_m=\cnorm{(I-P_m)d_{m+1}}_*>0 .
 \end{gathered}
\]
Define $\theta_m=|a_*(e_m,d_{m+1})|/S_m^{\rm fin}$ if
$S_m^{\rm fin}>0$, and $\theta_m=0$ otherwise.
The vector $(I-P_m)d_{m+1}$ is the innovation of the candidate
$d_{m+1}$, that is, its component orthogonal to $\widehat\Phi_m$, and
$\ell_m$ is the norm of that innovation.
Let $Q_0=I-P_{m_0}$ and assume explicit coefficients satisfy
\begin{equation}
 v_B=\sum_{d\in\cD_{\rm fin}}a_dQ_0d,\quad
 \sum_d|a_d|\leq B_{\rm entry},\quad
 \cnorm{Q_0u^*-v_B}_*\leq\sigma,\qquad B_{\rm entry}\geq0.
 \label{eq:c5-entry-source}
\end{equation}
The vectors $Q_0d$ are not renormalized, and $B_{\rm entry}$
need not equal the global atomic radius $B_{\rm src}$.

\begin{theorem}[Weighted finite-trajectory source bound]
\label{thm:c5-source}
Under \eqref{eq:c5-spaces}--\eqref{eq:c5-entry-source}, take
$0\leq\underline\theta_m\leq\theta_m$ and set
\[
 W_j=\sum_{m=m_0}^{m_0+j-1}\frac{\underline\theta_m^2}{\ell_m^2},
 \qquad x_0=(r_{m_0}^2-\sigma^2)_+ .
\]
If $B_{\rm entry}>0$ and $x_0>0$, then for $0\leq j\leq J$,
\begin{equation}
 (r_{m_0+j}^2-\sigma^2)_+
 \leq\left(x_0^{-1}+\frac{W_j}{4B_{\rm entry}^2}\right)^{-1}.
 \label{eq:c5-source-initial}
\end{equation}
In particular, if $W_j>0$,
$r_{m_0+j}^2\leq\sigma^2+4B_{\rm entry}^2/W_j$.
If $x_0=0$ or $B_{\rm entry}=0$, all later errors are at most $\sigma$.
For $W_j=0$ use projection monotonicity rather than division by zero.
\end{theorem}
\begin{proof}
For $w_m=(I-P_m)d_{m+1}$, Pythagoras gives
\begin{equation}
 r_m^2-r_{m+1}^2
 =\frac{|a_*(e_m,d_{m+1})|^2}{\ell_m^2}.
 \label{eq:c5-drop}
\end{equation}
The residual is orthogonal to both the starting and current spaces.
Consequently
\[
 r_m^2=a_*(e_m,Q_0u^*)\leq B_{\rm entry}S_m^{\rm fin}+\sigma r_m,
 \qquad
 S_m^{\rm fin}\geq\frac{(r_m^2-\sigma^2)_+}{2B_{\rm entry}}
\]
when $B_{\rm entry}>0$; the last step uses
$2(r_m^2-\sigma r_m)\geq r_m^2-\sigma^2$.
While both successive errors exceed $\sigma$, put
$x_m=r_m^2-\sigma^2>0$. Equation~\eqref{eq:c5-drop} implies
$x_{m+1}\leq x_m-b_mx_m^2$ with
$b_m=\underline\theta_m^2/(4B_{\rm entry}^2\ell_m^2)$.
Taking reciprocals proves \eqref{eq:c5-source-initial} up to a crossing.
After a crossing, monotonicity makes the positive part zero.
If $S_m^{\rm fin}=0$ or $B_{\rm entry}=0$, the source pairing gives
$r_m^2\leq\sigma r_m$ directly. This also covers zero errors.
\end{proof}

Barron--Cohen--Dahmen--DeVore prove in
\cite[Theorem~2.3, equation~(2.31)]{BarronCohenDahmenDeVore2008}
that exact OGA in a Hilbert space satisfies
\begin{equation}
 \cnorm{f-f_N}_H^2\leq\cnorm{f-h}_H^2+
            \frac{4\cnorm h_{L^1(\cD)}^2}{N},\qquad h\in L^1(\cD).
 \label{eq:c5-barron}
\end{equation}
Here $L^1(\cD)$ denotes the atomic variation space, not a spatial
Lebesgue norm. The correspondence in our restarted setting is
$f=Q_0u^*$, $h=v_B$, $\cnorm{f-h}_*\leq\sigma$, and atomic budget
at most $B_{\rm entry}$ on the unrenormalized projected dictionary.
The source pairing and reciprocal step above are the same classical
mechanism. Theorem~\ref{thm:c5-source} retains the observed
$\underline\theta_m^2/\ell_m^2$ along spaces selected by a nonlinear
algorithm. Neither the approximate-source remainder nor the projection
identity is claimed as a new general OGA principle.
If $\underline\theta_m\geq\underline\theta>0$, then $\ell_m\leq1$
gives $W_j\geq j\underline\theta^2$.
Stronger decay of the displayed bound requires independently controlled
growth of $W_j$ or stronger drop information.

For a direct comparison, Li--Siegel, Corollary~4.2 and equation~(4.11),
already use an approximate source \(h\) through an inequality of the form
\[
 \frac12\big(\|r\|^2-\|f-h\|^2\big)
 \leq \|h\|_{\mathcal A}\,|\langle r,g\rangle|,
\]
which leads to the same innovation-weighted projection recursion. The
present theorem specializes this mechanism to the restarted, unrenormalized
projected dictionary and retains the observed
\(\underline\theta_m^2/\ell_m^2\) weights generated by the nonlinear selection rule. It
should therefore be read as a specialization to the computed sequence of approximation
spaces, rather than as a new general approximate-source principle.

\subsection{An exact quartic energy connection}
\label{sec:c5-quartic}

For the one-dimensional pure or regularized quartic energy define
\[
 X=\{u\in W^{1,4}(0,1):\textstyle\int_0^1u=0\},\qquad
 \cE_\epsilon(u)=\int_0^1
       \left(\tfrac{\epsilon^2}2|u'|^2+\tfrac14|u'|^4\right)\cdd x-\ell(u).
\]
An additive constant in the density is immaterial.
Assume $u^*\in X$ solves the weak equation and let $Y\subset X$
contain $u^*$ and the spaces/dictionary of
Section~\ref{sec:c5-source}. Use
\[
 a_*(v,w)=\int_0^1(\epsilon^2+3|(u^*)'|^2)v'w'\cdd x .
\]
For $\epsilon>0$ this is positive definite on $Y$.
For $\epsilon=0$ assume explicitly that
$\int |(u^*)'|^2|v'|^2>0$ for every nonzero $v\in Y$;
the sufficient condition $(u^*)'\ne0$ almost everywhere ensures this
on the zero-mean space. Finiteness follows from Holder's inequality.
No full-space $H^1$ coercivity is inferred from positivity on $Y$.

For any $\widehat G_m\in\widehat\Phi_m$, define
\[
 D_m=\cnorm{\widehat G_m-z_m}_*,\quad
 v_m=\widehat G_m'-(u^*)',\quad
 R_m=\int_0^1\left((u^*)'v_m^3+\tfrac14v_m^4\right)\cdd x .
\]
\begin{proposition}[Quartic identity and comparison]
\label{prop:c5-quartic}
For every such state, with no exact-minimization requirement,
\begin{equation}
 \widehat\delta_m=\tfrac12(r_m^2+D_m^2)+R_m.
 \label{eq:c5-quartic}
\end{equation}
If $D_m\leq q_{\rm corr}r_m$ and
$|R_m|\leq\rho(r_m^2+D_m^2)$ throughout a finite range, where
$q_{\rm corr}\geq0$ and $0\leq\rho<1/2$, then
\begin{equation}
 (\tfrac12-\rho)r_m^2\leq\widehat\delta_m
 \leq C_Er_m^2,\qquad C_E=(\tfrac12+\rho)(1+q_{\rm corr}^2).
 \label{eq:c5-quartic-comparison}
\end{equation}
If $U_m$ is any source upper bound for $r_m^2$, the unconditional
additive estimate is
$\widehat\delta_m\leq U_m/2+D_m^2/2+|R_m|$.
\end{proposition}
\begin{proof}
Set $a=(u^*)'$ and $v=v_m$. Subtracting the linear part of the density
at $a$ gives
$\epsilon^2v^2/2+3a^2v^2/2+av^3+v^4/4$.
The weak equation cancels the integrated linear part, including the load.
Also $\widehat G_m-z_m$ lies in the active space and
$u^*-z_m$ is orthogonal to it, so
$\cnorm{\widehat G_m-u^*}_*^2=D_m^2+r_m^2$.
This proves the identity; the bounds follow by substitution.
\end{proof}

The nonlinear correction and remainder are already evaluated at the computed
state in \eqref{eq:c5-quartic}. Adding a further optimization-error term
to that identity would count the same effect twice.
Quadrature error appears instead as a separately reported numerical defect.
The constants $q_{\rm corr}$ and $\rho$ computed using $u^*$ are
quantities evaluated using the manufactured solution; the identity alone does not provide
solution-independent estimates of them. A small $\sigma$ at the starting index need
not be small relative to a terminal projection error.

\subsection{Projection bounds under a power-type decrease condition}
\label{sec:c5-power}

Retain the projections of Section~\ref{sec:c5-source} and set
\[
 x_m=r_m^2,\qquad
 \gamma_m=\frac{|a_*(e_m,d_{m+1})|}{\ell_m}.
\]
Then \eqref{eq:c5-drop} is $x_{m+1}=x_m-\gamma_m^2$.
When $S_m^{\rm fin}>0$, $\gamma_m=\theta_mS_m^{\rm fin}/\ell_m$;
otherwise $\gamma_m=0$ by its direct definition.
The word fractional in the following condition refers only to the power
imposed on the residual, not to a fractional differential operator.

\begin{theorem}[Explicit power-type finite-range bound]
\label{thm:c5-power}
Let $\alpha>0$, $c_0>0$, and assume, at every step with $r_m>0$
in $m_0\leq m<m_0+J$,
\begin{equation}
 \gamma_m^2\geq c_0r_m^{2+2/\alpha}.
 \label{eq:c5-power-condition}
\end{equation}
For $r_{m_0}>0$ and $0\leq j\leq J$,
\begin{equation}
 r_{m_0+j}^2\leq
 \left(r_{m_0}^{-2/\alpha}+\frac{c_0j}{\alpha}\right)^{-\alpha}.
 \label{eq:c5-power}
\end{equation}
If $r_{m_0}=0$, every later projection error is zero.
For the quartic models, if the relative conditions of
Proposition~\ref{prop:c5-quartic} also hold, then
\begin{equation}
 \widehat\delta_{m_0+j}\leq C_E
 \left(r_{m_0}^{-2/\alpha}+\frac{c_0j}{\alpha}\right)^{-\alpha}.
 \label{eq:c5-power-energy}
\end{equation}
Natural-distance and Sobolev bounds are obtained by taking respectively
the one-half and one-quarter powers, with the constants of
\CGAEnergyTheoremReference.
\end{theorem}
\begin{proof}
For $x_{m+1}>0$, convexity gives
\[
 x_{m+1}^{-1/\alpha}-x_m^{-1/\alpha}
 \geq\frac{x_m-x_{m+1}}{\alpha x_m^{1+1/\alpha}}
 \geq\frac{c_0}{\alpha}.
\]
Summation proves \eqref{eq:c5-power}. If a residual reaches zero,
nestedness gives the remaining cases. Substitute the projection bound
into \eqref{eq:c5-quartic-comparison}, then apply
\CGAEnergyTheoremReference.
\end{proof}

A sufficient factorization on positive-error steps is
\begin{equation}
 S_m^{\rm fin}\geq c_1r_m^{2\nu},\qquad
 \theta_m/\ell_m\geq c_2,\qquad
 \nu=\tfrac12+\tfrac1{2\alpha},\qquad c_1,c_2>0 .
 \label{eq:c5-power-factors}
\end{equation}
Indeed $\gamma_m^2\geq(c_1c_2)^2r_m^{4\nu}$ and
$4\nu=2+2/\alpha$. Thus $c_0=(c_1c_2)^2$ is admissible.
The residual visibility factor and the selected-direction factor are
separate requirements.

For a finite sequence with strictly positive drops, every prescribed
$\alpha>0$ admits
\[
 c_{\min}(\alpha;J)=
 \min_{\substack{m_0\leq m<m_0+J\\x_m>0}}
 \frac{x_m-x_{m+1}}{x_m^{1+1/\alpha}}>0 .
\]
Using this value on the same sequence constructs an a posteriori
envelope. It neither identifies a unique exponent nor independently
guarantees a rate for CGA. Evaluating a fresh finite range with the constants
fixed in advance supplies additional numerical evidence, but it still
does not prove an asymptotic theorem.
If any positive-error step has zero drop, no positive $c_0$ satisfies
\eqref{eq:c5-power-condition} on that whole finite range.
For a fixed finite dictionary whose span does not contain the target,
the positive best-approximation floor also prevents an indefinitely
valid positive $c_0$. Explicit bounds, their constants, and their
finite-range lengths are therefore more informative than unqualified
high-order notation.

\ifdefined\CGAChapterFiveEmbedded
\else
\begin{thebibliography}{9}
\bibitem{BarronCohenDahmenDeVore2008}
A.~R. Barron, A. Cohen, W. Dahmen, and R.~A. DeVore,
Approximation and learning by greedy algorithms,
\emph{Annals of Statistics} 36 (2008), 64--94.
\href{https://doi.org/10.1214/009053607000000631}{doi:10.1214/009053607000000631}.
\bibitem{LiSiegelEntropy2024}
Y. Li and J.~W. Siegel, Entropy-based convergence rates of greedy algorithms,
\emph{Mathematical Models and Methods in Applied Sciences} 34 (2024), 779--802.
\href{https://doi.org/10.1142/S0218202524500143}{doi:10.1142/S0218202524500143}.
The proof numbering used here is from
\href{https://arxiv.org/abs/2304.13332v2}{arXiv:2304.13332v2},
23 February 2024 (typeset 26 February).
\bibitem{DereventsovTemlyakov2022}
A. Dereventsov and V.~N. Temlyakov,
Biorthogonal greedy algorithms in convex optimization,
\emph{Applied and Computational Harmonic Analysis} 60 (2022), 489--511.
\href{https://doi.org/10.1016/j.acha.2022.05.001}{doi:10.1016/j.acha.2022.05.001}.
\bibitem{XuXu2025}
J. Xu and X. Xu,
Randomized greedy algorithms for neural network optimization in solving
partial differential equations,
\emph{Journal of Scientific Computing} 105 (2025), article 26.
\href{https://doi.org/10.1007/s10915-025-03050-5}{doi:10.1007/s10915-025-03050-5}.
\end{thebibliography}
\fi
\ifdefined\CGAChapterFiveStandalone
\def\CGAChapterFiveEndDocument{\end{document}}
\else
\let\CGAChapterFiveEndDocument\relax
\fi
\CGAChapterFiveEndDocument
